\section*{Bab \ref{ch:g4:air-a-mixture-of-gases} --- Udara, Campuran Gas}

\begin{solution}{exo:g4:air-a-mixture-of-gases:1}
Campuran. Dua gas utamanya adalah nitrogen dan oksigen.
\end{solution}

\begin{solution}{exo:g4:air-a-mixture-of-gases:2}
Sekitar empat perlima \omterm{def:g4:air-a-mixture-of-gases:air}{udara} adalah nitrogen, sekitar seperlima adalah
oksigen.
\end{solution}

\begin{solution}{exo:g4:air-a-mixture-of-gases:3}
Sekitar 78 liter nitrogen dan sekitar 21 liter oksigen.
\end{solution}

\begin{solution}{exo:g4:air-a-mixture-of-gases:4}
Botol itu tidak kosong: botol itu penuh \omterm{def:g4:air-a-mixture-of-gases:air}{udara}, dan \omterm{def:g4:air-a-mixture-of-gases:air}{udara} yang
terperangkap menahan air agar tidak masuk.
\end{solution}

\begin{solution}{exo:g4:air-a-mixture-of-gases:5}
Oksigennya lebih sedikit (16 bagian per 100, bukan 21) dan karbon
dioksidanya lebih banyak (4 bagian per 100, bukan hampir nol).
\end{solution}

\begin{solution}{exo:g4:air-a-mixture-of-gases:6}
Oksigen dalam kedua kasus: nyala api dan tubuh kita sama-sama memakai
oksigen dari \omterm{def:g4:air-a-mixture-of-gases:air}{udara}.
\end{solution}

\begin{solution}{exo:g4:air-a-mixture-of-gases:7}
78 kotak berupa nitrogen; $100 - 78 = 22$ kotak bukan nitrogen.
\end{solution}

\begin{solution}{exo:g4:air-a-mixture-of-gases:8}
Nyala api memakai oksigen dari \omterm{def:g4:air-a-mixture-of-gases:air}{udara} di dalam stoples. Ketika oksigen
di sekitar nyala tinggal terlalu sedikit, nyala itu padam.
\end{solution}

\begin{solution}{exo:g4:air-a-mixture-of-gases:9}
Seperlima dari 500 liter: $500 \div 5 = 100$ liter oksigen.
\end{solution}

\begin{solution}{exo:g4:air-a-mixture-of-gases:10}
Orang-orang di ruangan itu menghirup oksigen dan mengembuskan karbon
dioksida. Membuka jendela memasukkan kembali \omterm{def:g4:air-a-mixture-of-gases:air}{udara} segar, yang lebih
banyak oksigennya dan lebih sedikit karbon dioksidanya.
\end{solution}

\begin{solution}{exo:g4:air-a-mixture-of-gases:11}
\omterm{def:g4:air-a-mixture-of-gases:air}{Udara} tiga kali lebih banyak mengandung oksigen tiga kali lebih banyak:
kira-kira $3 \times 10 = 30$ detik.
\end{solution}
